\section{Introduction}
\label{sec:intro}

Deep neural networks have achieved remarkable success across computer vision tasks, yet our understanding of their internal representations remains limited. While significant progress has been made in analyzing activation spaces~\cite{carlssonTopologicalApproachesDeep2018}, the geometry and topology of \emph{weight spaces}---the high-dimensional manifolds where network parameters reside---remain largely unexplored. Understanding weight space structure is crucial for addressing fundamental questions: Why do different optimization trajectories converge to similar solutions? How do architectural choices affect the loss landscape? Can we predict generalization from weight space geometry?

We introduce \textbf{hyper-representations}, a novel framework that treats neural network weights as data points to be embedded and analyzed. Unlike traditional approaches that study individual networks in isolation, hyper-representations enable comparative analysis across models, architectures, and training regimes. By learning compact embeddings of weight vectors, we can visualize, cluster, and reason about entire populations of neural networks.

Our key insight is that \emph{multiparameter persistent homology}~\cite{carlssonTheoryMultidimensionalPersistence2009} provides a principled mathematical framework for extracting robust topological features from weight space trajectories. Traditional single-parameter persistence tracks how topological features evolve along one filtration (e.g., scale). However, neural network training involves multiple interacting parameters: learning rate, regularization strength, data overlap, and loss function choice. Multiparameter persistence~\cite{botnanBottleneckStabilityRank2024} naturally captures these multi-scale phenomena, revealing topological signatures invisible to conventional analysis.

\subsection{Contributions}

We make the following contributions:

\begin{enumerate}
    \item \textbf{Hyper-representation framework}: We introduce a transformer-based autoencoder architecture that learns compact, semantically meaningful embeddings of neural network weights, preserving both geometric and topological structure.
    
    \item \textbf{Multiparameter persistence analysis}: We apply multiparameter persistent homology to weight space trajectories, revealing how different loss functions induce distinct topological signatures. Our analysis identifies stable topological features that correlate with generalization performance.
    
    \item \textbf{Systematic loss function study}: We evaluate 18 diverse loss functions spanning reconstruction losses (MSE, MAE), spectral losses (FFT, MelSpec), optimal transport (Sinkhorn), information-theoretic measures (KL, JS), and topological losses (persistence-based). We demonstrate that regularized combinations exhibit superior stability and generalization.
    
    \item \textbf{Statistical and topological insights}: Through comprehensive analysis of 51 trained models across three data overlap scenarios, we establish strong correlations between topological features (effective rank, persistence entropy, Betti numbers) and model performance.
\end{enumerate}

Our work bridges topological data analysis~\cite{ballesterTopologicalDataAnalysis2024} and deep learning, demonstrating that weight space topology is not merely a theoretical curiosity but a practical tool for understanding and improving neural networks. The hyper-representation framework enables new applications: model selection via topological clustering, transfer learning through weight space interpolation, and architecture search guided by topological constraints.

\subsection{Why Multiple Loss Functions Matter}

The choice of loss function profoundly affects the optimization landscape and learned representations. We systematically study diverse loss families:

\begin{itemize}
    \item \textbf{Reconstruction losses} (MSE, MAE, MAPE): Establish baseline weight space geometry through direct distance minimization.
    \item \textbf{Spectral losses} (FFT, MelSpec): Capture frequency-domain structure, revealing periodic patterns in weight evolution.
    \item \textbf{Optimal transport} (Sinkhorn): Preserve distributional properties, inducing smoother weight space trajectories.
    \item \textbf{Information-theoretic} (KL, JS): Enforce probabilistic structure, promoting diverse representations.
    \item \textbf{Geometric losses} (Frobenius, LogNorm): Control weight matrix norms and conditioning.
    \item \textbf{Topological losses} (Persistence-based): Directly optimize for desired topological features.
    \item \textbf{Regularized combinations}: Combine complementary losses for improved stability and generalization.
\end{itemize}

Each loss family explores different regions of weight space, and their regularized combinations create rich, multi-scale optimization landscapes. Our multiparameter persistence analysis reveals that these combinations exhibit more stable topological features than individual losses, explaining their superior generalization~\cite{birdalIntrinsicDimensionPersistent2021}.

The remainder of this paper is organized as follows: Section~\ref{sec:related} reviews related work in topological data analysis and neural network analysis. Section~\ref{sec:method} details our hyper-representation framework and multiparameter persistence methodology. Section~\ref{sec:experiments} describes our experimental setup and loss function taxonomy. Section~\ref{sec:results} presents comprehensive statistical and topological analysis. Section~\ref{sec:conclusion} concludes with implications and future directions.

\subsection{Mathematics}

Please, number all of your sections and displayed equations as in these examples:
\begin{equation}
  E = m\cdot c^2
  \label{eq:important}
\end{equation}
and
\begin{equation}
  v = a\cdot t.
  \label{eq:also-important}
\end{equation}
It is important for the reader to be able to refer to any particular equation.
Just because you did not refer to it in the text does not mean that some future reader might not need to refer to it.
It is cumbersome to have to use circumlocutions like ``the equation second from the top of page 3 column 1''.
(Note that the ruler will not be present in the final copy, so is not an alternative to equation numbers).
All authors will benefit from reading Mermin's description of how to write mathematics:
\url{http://www.pamitc.org/documents/mermin.pdf}.
